\documentclass[aps,pra,reprint,amsmath,amssymb,floatfix,nofootinbib]{revtex4-2}
\usepackage{amsthm,mathtools,mathrsfs,booktabs}
\usepackage[T1]{fontenc}
\usepackage{newtxtext,newtxmath}
\usepackage[colorlinks=true,linkcolor=blue,citecolor=blue,urlcolor=blue]{hyperref}
\usepackage[nameinlink,noabbrev]{cleveref}
\newcommand{\Tr}{\operatorname{Tr}}
\newcommand{\rk}{\operatorname{rank}}
\newcommand{\ran}{\operatorname{ran}}
\newcommand{\End}{\operatorname{End}}
\newcommand{\vecop}{\operatorname{vec}}
\newcommand{\Sym}{\operatorname{Sym}}
\newcommand{\norm}[1]{\left\lVert#1\right\rVert}
\newcommand{\ip}[2]{\left\langle#1,#2\right\rangle}
\newcommand{\C}{\mathbb C}
\newcommand{\R}{\mathbb R}
\newcommand{\Q}{\mathbb Q}
\newcommand{\Sfour}{\mathfrak S_4}
\newcommand{\Sthree}{\mathfrak S_3}
\newtheorem{theorem}{Theorem}
\newtheorem{lemma}[theorem]{Lemma}
\newtheorem{proposition}[theorem]{Proposition}
\newtheorem{corollary}[theorem]{Corollary}
\setlength{\emergencystretch}{2em}
\begin{document}
\title{Three-copy undistillability of the boundary two-ququart Werner state}
\author{Ao-Xiang Liu}
\noaffiliation
\date{15 October 2026}
\begin{abstract}
We give a computer-assisted proof that the two-ququart Werner state
$\rho_*=(I-F/2)/14$ is three-copy undistillable.
Equivalently, its three-factor partial-trace form is nonnegative on every
complex $64\times64$ matrix of rank at most two, including nonnormal
matrices with unequal singular values.
An isometric factorization embeds every such matrix into three matched
four-replica moments with valid partial-transpose and joint positivity
constraints. Common symmetry averaging gives a complete rational source
space of dimension $2741$, including the additional alternating component
present at local dimension four.
An exact functional identity combines $1239$ rational positive Gram
multipliers with homogeneous equality constraints to reproduce the
objective. Every coefficient and positivity inequality is verified in
rational arithmetic, with a second implementation rebuilding the
coefficient maps from literal permutation matrices.
The result gives the exact three-copy distillability threshold
$\alpha<-1/2$ for the two-ququart Werner family.
\end{abstract}
\maketitle

\section{Introduction}
\label{sec:introduction}

Entanglement distillation distinguishes quantum correlations that can
supply pure entangled pairs from correlations that remain inaccessible
to local operations and classical communication. States with positive
partial transpose are undistillable, even when arbitrarily many copies
are available~\cite{HHH}. Whether this obstruction exhausts bound
entanglement is a longstanding question. A state with negative partial
transpose that remained undistillable at every copy number would establish
the existence of negative-partial-transpose bound entanglement.

Werner states provide a symmetric setting in which to study this question
\cite{Werner,DCLB,DSS,Five}. Their symmetry reduces the state family to one
parameter, while the distillation test still ranges over arbitrary
Schmidt-rank-two vectors across all available copies. The boundary of
one-copy distillability is particularly useful. Proving that this boundary
state remains undistillable at a larger fixed copy number isolates the
possibility that collective processing may activate its entanglement.
The two-ququart instance at two copies was singled out in
Ref.~\cite{Five}, following matrix formulations developed in
Ref.~\cite{PPHH}.

The finite-copy criterion admits a formulation through partial traces of
low-rank coefficient matrices~\cite{CostaRico}.
Recent preprints establish the sharp two-copy inequality and its Werner
consequences~\cite{SC,FHPV,FGP,BGH}. Extending that argument to three copies
requires new control of the simultaneous reductions of a matrix.
Partial trace can increase rank, and the two-copy inequality cannot be
applied successively while retaining the rank-two hypothesis.
Wu and Zou prove the three-copy endpoint for all normal rank-two matrices
and for specified nonnormal support classes in arbitrary local
dimensions~\cite{WZ}.
The remaining matrices need not have either of these structures.

Semidefinite formulations of Werner distillability have a substantial
history~\cite{VD}. A recent candidate manuscript gives a rational
certificate for the unrestricted three-copy qutrit endpoint using
matched four-replica moments and a coupled positive block~\cite{Candidate}.
We use that moment construction, with its unrestricted second matrix
factor, and extend the source coordinates and positive representations
to local dimension four. This extension introduces a fourth alternating
component and its partially transposed representations. A certificate in
the qutrit quotient does not control these additional directions.

Here we construct a rational certificate for the unrestricted two-ququart
endpoint. The proof separates the physical reduction from a finite
algebraic verification. Every complex matrix of rank at most two produces
feasible matched moments. A complete rational basis represents their
common symmetry average, and an exact positive-multiplier identity proves
nonnegativity throughout the resulting relaxation.
The certificate is accompanied by a second implementation that rebuilds
its source and Gram maps before checking the rational identity.
The two-copy theorems and the claimed qutrit endpoint are context for this
construction, rather than premises of its proof.

\section{Main result and the distillation criterion}
\label{sec:main}

On $\C^4\otimes\C^4$, the Werner family is
\begin{align}
 \rho(4,\alpha)=\frac{I+\alpha F}{16+4\alpha},
 \qquad -1\leq\alpha\leq1\ ,
 \label{eq:werner}
\end{align}
where $F|i,j\rangle=|j,i\rangle$.
The state of interest is $\rho_* = \rho(4,-1/2)$.
Let $H=H_1\otimes H_2\otimes H_3$, with $H_i=\C^4$.
For $C\in\End(H)$, partial trace contracts matching row and column indices
at the indicated sites. It is unnormalized, and $\Tr_\varnothing C=C$.
The relevant quadratic form is
\begin{align}
 q_3(C)
 &=\sum_{J\subseteq\{1,2,3\}}
       \left(-\frac12\right)^{|J|}\norm{\Tr_JC}_F^2\nonumber\\
 &=\norm C_F^2-\frac12\sum_{i=1}^3\norm{\Tr_iC}_F^2\nonumber\\
 &\quad+\frac14\sum_{i<j}\norm{\Tr_{ij}C}_F^2
   -\frac18|\Tr C|^2\ .
 \label{eq:q3}
\end{align}
All matrix spaces in this paper are complex unless a real or rational
coefficient space is explicitly specified. The Frobenius inner product
is $\ip XY=\Tr(X^\dagger Y)$.

\begin{theorem}[Three-copy endpoint]
\label{thm:d4-three-certificate}
Every $C\in M_{64}(\C)$ with $\rk C\leq2$ satisfies $q_3(C)\geq0$.
Consequently, $\rho_*=(I-F/2)/14$ is three-copy undistillable.
The proof is computer assisted and uses the exact rational certificate
specified in \cref{sec:certificate,app:verification}.
\end{theorem}

The theorem covers arbitrary input and output support planes and arbitrary
singular values. The inequality is sharp. For any unit
$u,v\in\C^4\otimes\C^4$, the matrix
\begin{align}
 C_0=\frac{|0\rangle\langle0|+|1\rangle\langle1|}{2}
                 \otimes |u\rangle\langle v|
 \label{eq:product-zero}
\end{align}
has rank two, $\norm{C_0}_F^2=1/2$, and $q_3(C_0)=0$.
Indeed, the partial-trace form factorizes across this tensor product,
and its first factor has one-copy value
$1/2-|1|^2/2=0$.
For generic $u,v$ the resulting matrix is nonnormal.

The connection to distillation follows directly from vectorization.
The vector associated with $C$ is
$|\psi_C\rangle=\sum_{x,y}C_{xy}|x\rangle_A|y\rangle_B$.
Its Schmidt rank across $A_1A_2A_3|B_1B_2B_3$ equals $\rk C$.
With $|\Omega\rangle=\sum_{j=0}^3|jj\rangle$, partial transposition
on Bob's factor gives
\begin{align}
 \rho_*^{\Gamma_B}
 =\frac1{14}\left(I-\frac12|\Omega\rangle\langle\Omega|\right)\ .
 \label{eq:werner-pt}
\end{align}
Contracting the $\Omega$ factors for a subset $J$ leaves
$\Tr_JC$ and its complex conjugate. Expansion therefore yields
\begin{align}
 \langle\psi_C|\bigl(\rho_*^{\Gamma_B}\bigr)^{\otimes3}
                  |\psi_C\rangle
 =\frac{q_3(C)}{14^3}\ .
 \label{eq:distillation-form}
\end{align}
The Schmidt-rank-two criterion for finite-copy distillation
\cite{HHH,DCLB} proves the stated consequence.
The normalized vector $\Omega/2$ has eigenvalue $-1/14$ in
\cref{eq:werner-pt}, confirming that $\rho_*$ has negative partial
transpose.

The endpoint also determines the full three-copy parameter threshold.

\begin{corollary}[Two-ququart threshold]
\label{cor:threshold}
For $-1\leq\alpha\leq1$, the state $\rho(4,\alpha)$ is three-copy
distillable if and only if $\alpha<-1/2$.
In particular every state with $-1/2\leq\alpha<-1/4$ is NPT and
three-copy undistillable.
\end{corollary}

\begin{proof}
For $-1/2\leq\alpha\leq0$, a direct calculation gives
\begin{align}
 \rho(4,\alpha)
 =p\rho_*+(1-p)\frac I{16},
 \qquad p=\frac{-7\alpha}{4+\alpha}\in[0,1]\ .
 \label{eq:local-replacement}
\end{align}
The parties can implement this map by shared randomness, retaining their
input with probability $p$ and otherwise locally preparing maximally
mixed states. Applying this local operation to each of three copies
cannot turn a three-copy undistillable state into a distillable one.
For $\alpha\geq0$, the partial transpose is positive.
For $\alpha<-1/2$, the vector
$(|00\rangle+|11\rangle)/\sqrt2$ gives the one-copy expectation
$(1+2\alpha)/(16+4\alpha)<0$.
Tensoring it with $|0\rangle_A|1\rangle_B$ on each remaining copy
preserves Schmidt rank two and gives a negative three-copy expectation.
Finally, the only potentially negative eigenvalue of
$\rho(4,\alpha)^{\Gamma_B}$ is
$(1+4\alpha)/(16+4\alpha)$, which proves the NPT statement.
\end{proof}

\subsection{Relation to the fourth and fifth problems}
\label{sec:problem-reductions}

Problem 4 of Ref.~\cite{Five} asks whether an NPT state can remain
undistillable at every finite copy number. For a fixed state $\rho$,
the distillation criterion is
\begin{align}
 \rho\text{ is distillable}
 \quad\Longleftrightarrow\quad
 &\exists n\geq1,\ \exists\psi,\quad \operatorname{SR}(\psi)\leq2,
 \nonumber\\[-2pt]
 &\langle\psi|(\rho^{\Gamma_B})^{\otimes n}|\psi\rangle<0\ .
 \label{eq:general-distillation-criterion}
\end{align}
Here Schmidt rank refers to the bipartition between all of Alice's
systems and all of Bob's systems. The two-dimensional local projections
can act jointly on all copies held by each party. Their output is a
two-qubit NPT state, which is distillable~\cite{HHH,DCLB}.
For NPT bound entanglement, the same state must fail this test for
every $n$. Establishing failure at one fixed $n$ does not supply that
quantifier.

The existence question can be restricted to Werner states without loss
of generality~\cite{HHWerner,DCLB,PPHH}.
For the nontrivial direction, suppose that $\rho$ is NPT and
undistillable. A vector $\psi$ has negative expectation against
$\rho^{\Gamma_B}$. After embedding the local spaces in equal dimension,
a local filter $L$ can be chosen with
$(L^\dagger\otimes I)\Omega_d$ proportional to $\psi$.
Its successful output $\sigma$ satisfies
$\Tr(F\sigma)=\langle\Omega_d|\sigma^{\Gamma_B}|\Omega_d\rangle<0$.
Averaging $\sigma$ under $U\otimes U$ produces a Werner state and
preserves this negative swap expectation, hence its NPT property.
Filtering and averaging are local operations with classical
communication. A distillation protocol for their output would also
distill the input, which proves the required preservation of
undistillability. The converse follows because Werner states are
bipartite states.

In the convention
$\rho(d,\alpha)=(I+\alpha F)/(d^2+\alpha d)$, NPT is equivalent to
$\alpha<-1/d$. One-copy distillability is equivalent to $\alpha<-1/2$.
Indeed, the squared overlap of a normalized Schmidt-rank-two vector
with the unnormalized $\Omega_d=\sum_{j=0}^{d-1}|jj\rangle$ is at most
two, and that value is attained. Problem 4 is therefore equivalent to
the existence of $d>2$ and $\alpha\in[-1/2,-1/d)$ such that
\begin{align}
 \langle\psi|[\rho(d,\alpha)^{\Gamma_B}]^{\otimes n}|\psi\rangle
 \geq0
 \label{eq:problem-four-quantifiers}
\end{align}
for every $n\geq1$ and every $\psi$ of Schmidt rank at most two.
The dimension and the state are chosen once, before the universal
quantifier over the copy number.

Problem 5 selects $d=4$, $\alpha=-1/2$, and $n=2$~\cite{Five}.
Theorem~\ref{thm:d4-three-certificate} concerns the same state at $n=3$.
Selecting this endpoint is a sufficient route to a positive solution of
Problem 4 if it works for every $n$. It is not an equivalent restriction
of its existence question. Failure of this particular route would not
exclude other NPT bound entangled Werner states.

The distinguished dimension four makes the partial transpose proportional
to a reflection. With $P=|\Omega\rangle\langle\Omega|/4$, its
unnormalized form is $I-2P$. For two copies,
\begin{align}
 14^2(\rho_*^{\Gamma_B})^{\otimes2}&=I-2Q_2,\nonumber\\
 Q_2&=P\otimes(I-P)+(I-P)\otimes P\ .
 \label{eq:half-property-projector}
\end{align}
Two-copy undistillability is equivalent to the half-property
$\langle\psi|Q_2|\psi\rangle\leq1/2$ for every normalized
Schmidt-rank-two vector~\cite{PPHH}.
Under vectorization, the normalized vectors in $\ran Q_2$ have coefficient
matrices $K=A\otimes I_4+I_4\otimes B$, where
$\Tr A=\Tr B=0$ and $\norm A_F^2+\norm B_F^2=1/4$.
Optimizing the squared overlap of a fixed vector with Schmidt-rank-two
vectors gives the sum of its two largest squared Schmidt coefficients.
Exchanging the two maximizations consequently gives the equivalent
matrix inequality
\begin{align}
 s_1(K)^2+s_2(K)^2\leq\frac12\ .
 \label{eq:two-copy-singular-values}
\end{align}
These are singular values of $K$, or square roots of eigenvalues of
$K^\dagger K$. Eigenvalues of $K$ alone do not suffice for nonnormal
matrices. This is the reduction in Eqs.~(50)--(52) of Ref.~\cite{PPHH}.

For any copy number, vectorization provides a second exact matrix
formulation. If $C\in M_{4^n}(\C)$, the same contraction argument as in
\cref{eq:distillation-form} gives
\begin{align}
 q_n(C)&=\sum_{J\subseteq\{1,\ldots,n\}}
          \left(-\frac12\right)^{|J|}\norm{\Tr_J C}_F^2,\nonumber\\
 \langle\psi_C|(\rho_*^{\Gamma_B})^{\otimes n}|\psi_C\rangle
 &=14^{-n}q_n(C)\ .
 \label{eq:general-partial-trace-form}
\end{align}
The rank of $C$ equals the Schmidt rank of $\psi_C$. Thus nonnegativity
of $q_n$ on all complex matrices of rank at most two is equivalent to
$n$-copy undistillability of $\rho_*$. No Hermiticity, normality, or
equality of singular values is assumed. The case $n=2$ gives Problem 5,
and $n=3$ gives the present theorem. The matched moments below provide an
outer relaxation of this exact matrix problem. Their positivity gives
a sufficient certificate and does not assert that every feasible moment
has a physical realization.

We now establish \cref{thm:d4-three-certificate}.
The proof has three components. An exact lifting produces valid positive
moments for every matrix in the theorem. Complete source coordinates
represent their symmetry average. A rational certificate expresses the
objective as a sum of positive pairings modulo valid equalities.

\section{Matched moments for arbitrary complex rank-two matrices}
\label{sec:matched-four}

\subsection{An isometric lifting}

Homogeneity allows a nonzero $C$ to be normalized by
$\norm C_F^2=1/2$.
An isometry $A\in M_{64,2}(\C)$ with $\ran C\subseteq\ran A$
and the matrix $B=2C^\dagger A$ then satisfy
\begin{align}
 C=\frac12AB^\dagger,
 \qquad A^\dagger A=I_2,
 \qquad \norm B_F^2=2\ .
 \label{eq:isometric-lift}
\end{align}
The first identity follows from $AA^\dagger C=C$.
For the last one,
$\norm B_F^2=4\Tr(C^\dagger AA^\dagger C)=4\norm C_F^2$.
A rank-one range can always be completed to the required two-dimensional
range. There is no condition $B^\dagger B=I_2$.
In particular the lifting preserves unequal singular values of $C$.

The unit vectors
$a=\vecop(A)/\sqrt2$ and $b=\vecop(B)/\sqrt2$ lie in
$\mathcal V=H\otimes Q$, where $Q=\C^2$ and $\dim\mathcal V=128$.
The reduced state of $a$ on $Q$ is $I_2/2$.
For two replicas of $\mathcal V$, let $F_i$ swap their $H_i$ factors
and let $F_Q$ swap their $Q$ factors.
The swap identity takes the form
\begin{align}
 \langle a\otimes b|F_Q\prod_{j\in J}F_j|a\otimes b\rangle
 =\norm{\Tr_JC}_F^2
 \qquad(J\subseteq\{1,2,3\})\ .
 \label{eq:matched-swap-subsets}
\end{align}
To check the conjugations, let $s,t$ index the traced sites and $u,v$
the remaining sites. Both sides equal
\begin{align}
 \frac14\sum_{s,t,u,v,\mu,\nu}
 A_{su,\mu}\overline{B_{sv,\mu}}
 \overline{A_{tu,\nu}}B_{tv,\nu}\ .
 \label{eq:swap-indices}
\end{align}
It follows that
\begin{align}
 q_3(C)=\langle a\otimes b|
     F_Q\prod_{i=1}^3\left(I-\frac12F_i\right)|a\otimes b\rangle\ .
 \label{eq:matched-objective}
\end{align}

\subsection{Moment equalities and positive maps}

The matched-moment construction uses four replicas of $\mathcal V$.
These replicas are auxiliary moments and are distinct from the three
resource copies in the distillation problem.
For $P_v=|v\rangle\langle v|$, the three moments are
\begin{align}
 Y_{31}&=P_a\otimes P_a\otimes P_a\otimes P_b,\nonumber\\
 Y_{22}&=P_a\otimes P_a\otimes P_b\otimes P_b,\nonumber\\
 Y_{13}&=P_b\otimes P_b\otimes P_b\otimes P_a\ .
 \label{eq:moments}
\end{align}
The order in $Y_{13}$ is part of the convention.
Each moment and every partial transpose on complete replicas are
positive semidefinite. This follows because each is a tensor product
of positive rank-one projectors, and transposition on a complete replica
replaces its vector by its complex conjugate.

The normalization of $a,b$ gives two matching marginals. Deleting
replica 3 from $Y_{31}$ and $Y_{22}$ leaves $P_a\otimes P_a\otimes P_b$.
Deleting replica 2 from $Y_{22}$ leaves $P_a\otimes P_b\otimes P_b$.
Deleting replica 2 from $Y_{13}$ and ordering its remaining replicas as
$(4,1,3)$ gives the same operator.
For a replica $r$ occupied by $a$, the conditional equality is
\begin{align}
 \Tr_{H_r}Y_f
 =\frac{I_{Q_r}}2\otimes\Tr_{H_rQ_r}Y_f\ .
 \label{eq:matched-conditional}
\end{align}
The tensor factors on the right are restored to their original positions.
The construction imposes this equality at $r=1$ for $f=31,22$ and at
$r=4$ for $f=13$. The identical-factor symmetries imply the analogous
equalities at the other $a$ replicas. No such condition is imposed at a
$b$ replica.

A coupled positive block retains information shared by the three moments
\cite{Candidate}. The vectors
$z=a\otimes a\otimes b\otimes\bar a$ and
$w=b\otimes b\otimes a\otimes\bar b$ give
\begin{align}
 \mathcal J(Y)=
 \begin{pmatrix}zz^\dagger&zw^\dagger\\wz^\dagger&ww^\dagger\end{pmatrix}
 \succeq0\ .
 \label{eq:matched-joint}
\end{align}
All four entries are linear images of the moments.
The convention
$V_p|x_1,x_2,x_3,x_4\rangle
=|x_{p^{-1}(1)},x_{p^{-1}(2)},x_{p^{-1}(3)},x_{p^{-1}(4)}\rangle$
fixes the permutations $\tau=(3\ 4)$ and $\pi=(1\ 3\ 2\ 4)$.
The linear maps are
\begin{align}
 \mathcal D(Y)&=(V_\tau YV_\tau^{\mathsf T})^{\Gamma_4},
 &\mathcal K(Y)&=(YV_\pi)^{\Gamma_4}\ ,
 \label{eq:joint-maps}
\end{align}
and the block is
\begin{align}
 \mathcal J(Y)=
 \begin{pmatrix}
 \mathcal D(Y_{31})&\mathcal K(Y_{22})\\
 \mathcal K(Y_{22})^\dagger&\mathcal D(Y_{13})
 \end{pmatrix}\ .
 \label{eq:joint-linear}
\end{align}
For example, with arbitrary row and column multi-indices $x,y$ in
$\mathcal V^{\otimes4}$, the upper right entry has matrix element
\begin{align}
 a_{x_1}a_{x_2}b_{x_3}\bar a_{x_4}
 \bar b_{y_1}\bar b_{y_2}\bar a_{y_3}b_{y_4}
 =z_x\bar w_y\ .
 \label{eq:joint-cross-indices}
\end{align}
This follows by right multiplication by $V_\pi$ followed by transposition
on replica 4. The diagonal identities follow by exchanging replicas 3
and 4 before that transposition.
Thus the block inequality holds for general complex $a,b$.

\subsection{Common averaging}

We apply the same averaging operations to all three moments.
For $U=U_1\otimes U_2\otimes U_3\otimes U_Q$, with $U_i\in U(4)$
and $U_Q\in U(2)$, average $Y_f$ under $U^{\otimes4}$.
Also average over permutations of the three physical sites and over
complex conjugation. The resulting moments are denoted $\overline Y_f$.
Every equality above is preserved. For a complete-replica transpose,
unitary conjugation becomes conjugation by the corresponding tensor
product of $U$ and $\bar U$, which preserves positivity.

The coupled block is preserved as a whole. Under $a\mapsto Ua$ and
$b\mapsto Ub$, both $z$ and $w$ transform by
$U^{\otimes3}\otimes\bar U$. Physical-site permutations likewise act
identically on the two components. Complex conjugation sends the whole
block to its complex conjugate. The average is therefore positive.
Its entries remain the same linear images of the averaged moments.
This common averaging is necessary to preserve the cross entry.

The objective is unchanged. It is real and invariant under all these
operations. In the four-replica notation it is
$\Tr(OY_{31})$, where
\begin{align}
 O=F_Q^{(1,4)}\prod_{i=1}^3\left(I-\frac12F_i^{(1,4)}\right)\ .
 \label{eq:four-replica-objective}
\end{align}
The two unused replicas have unit trace.
Consequently every original matrix gives averaged moments with objective
$q_3(C)$, valid homogeneous equalities, and valid positive blocks.
This is an outer relaxation of the physical moment set.
No converse realization of arbitrary feasible moments is needed.
In particular, using real symmetric averaged moments imposes no reality
condition on $C$.

\section{Complete rational coordinates and positive compressions}
\label{sec:coordinates}

\subsection{Source generators and their normalization}

Regroup the four replicas by physical site and auxiliary qubit.
Schur--Weyl duality places the local unitary commutant in the span of
\begin{align}
 D_{\boldsymbol p}
 =V_{p_1}^{(4)}\otimes V_{p_2}^{(4)}\otimes V_{p_3}^{(4)}
                         \otimes V_{p_Q}^{(2)},
 \qquad \boldsymbol p\in\Sfour^4\ .
 \label{eq:diagrams}
\end{align}
The 24 permutations at $d=4$ are independent.
The auxiliary permutation span at $d=2$ has dimension 14.

Let $R_{31}=R_{13}$ be the subgroup of $\Sfour$ fixing replica 4,
and let $R_{22}$ preserve the two sets $\{1,2\}$ and $\{3,4\}$.
The global identical-factor projector is
$\Pi_f=|R_f|^{-1}\sum_{r\in R_f}V_r^{(\mathcal V)}$.
The moment $\overline Y_f$ is supported on its range.
For $W_\sigma$ permuting the physical sites simultaneously in all replicas,
the source projection is
\begin{align}
 \mathcal P_f(X)=\frac1{12}\sum_{\sigma\in\Sthree}
 W_\sigma\bigl[\Pi_fX\Pi_f+(\Pi_fX\Pi_f)^{\mathsf T}\bigr]
 W_\sigma^{\mathsf T}\ .
 \label{eq:source-projection}
\end{align}
It is an idempotent real linear map on the local commutant.
The source generators are $B_{f,k}=\mathcal P_f(D_{\boldsymbol p_{f,k}})$.
The chosen words $\boldsymbol p_{f,k}$ are recorded in the certificate.

\begin{lemma}[Complete source spaces]
\label{lem:source-complete}
The selected generators form rational bases of the source spaces for
$\overline Y_{31},\overline Y_{22},\overline Y_{13}$.
Their dimensions are $667,1407,667$, respectively.
\end{lemma}

\begin{proof}
The symmetry properties place each averaged moment in
$\ran\mathcal P_f$. The trace of this idempotent is its rank.
The exact trace calculation in \cref{app:source} gives $667$ for
$f=31,13$ and $1407$ for $f=22$.
Reducing the selected generators in a literal local permutation basis
produces integer coefficient matrices after columnwise denominator
clearing. Their square minors have nonzero determinants $686863$,
$137174$, and $686863$ modulo $1000033$.
These determinants prove independence over $\Q$ and $\R$.
The independence lower bounds equal the source dimension upper bounds,
which proves spanning.
\end{proof}

The stored coordinate convention includes a positive scale
$s=128^4=268435456$.
For $x=(x_{31},x_{22},x_{13})\in\R^{2741}$, define
$\mathcal Y_f(x)=\sum_k x_{f,k}B_{f,k}$.
The coordinates of physical moments obey
\begin{align}
 \mathcal Y_f(x)=s\overline Y_f\ .
 \label{eq:matched-source-scale}
\end{align}
The stored objective has
$c_{31,k}=s^{-1}\Tr(OB_{31,k})$ and $c_{22,k}=c_{13,k}=0$.
It follows that $c^{\mathsf T}x=q_3(C)$ exactly.
Equality and positive maps act on $\mathcal Y_f(x)$.
Their common factor $s$ preserves the homogeneous equalities and
positive semidefiniteness.

Expanding the conditional and matching maps in independent output
coordinates gives a $2912\times2741$ integer matrix $E$ after a common
denominator $1152$ is cleared. Every physical coordinate vector satisfies
$Ex=0$. Its exact rational rank is 272.
The certificate uses the submatrix $E_0$ comprising the 272 selected rows
listed in its equality file. Their rank is also 272.
Explicit partial-trace and ordering conventions are given in
\cref{app:equalities}.

At dimension four, the source space includes an additional local
alternating component. The antisymmetrizer
$\mathcal A_4^{(d)}=(1/24)\sum_{p\in\Sfour}\operatorname{sgn}(p)V_p^{(d)}$
is zero at $d=3$ and has rank one at $d=4$.
Its marginal is
\begin{align}
 \Tr_r\mathcal A_4^{(d)}=\frac{d-3}{4}\mathcal A_3^{(d)}\ .
 \label{eq:alternating-marginal}
\end{align}
Support on the alternating subspace and unitary invariance imply
proportionality, and traces give the factor
$\binom d4/\binom d3$.
The dimension-four construction retains this direction and its partially
transposed representations. In particular the source calculation uses
all 24 physical permutations rather than the 23-dimensional qutrit
quotient.

\subsection{Literal representations and Gram maps}

For each used local dimension $d$ and replica cut $T$, the archive supplies
rational embeddings $J_\lambda$ and local matrices $\rho_\lambda(p)$
satisfying the literal identities
\begin{align}
 (V_p^{(d)})^{\Gamma_T}J_\lambda
   =J_\lambda\rho_\lambda(p),
 \qquad G_\lambda=J_\lambda^{\mathsf T}J_\lambda\succ0\ .
 \label{eq:local-intertwiner}
\end{align}
Here $J_\lambda$ has full column rank.
All 24 identities are verified for each of the 50 local embeddings.
Tensoring four such embeddings gives $L_t$ and
$G_t=L_t^{\mathsf T}L_t$ for each parent block.
For a source block obtained from $\mathcal Y_f(x)^{\Gamma_T}$,
the rational matrix $\mathscr R_t(x)$ then obeys
\begin{align}
 L_t^{\mathsf T}\mathcal Y_f(x)^{\Gamma_T}L_t
   =G_t\mathscr R_t(x)\ .
 \label{eq:literal-compression}
\end{align}
The right side is symmetric, even when $\mathscr R_t(x)$ is not
symmetric in its nonorthonormal coordinates.

For each retained block the archive gives a rational frame $N_t$.
Its positive map is
\begin{align}
 \widehat M_t(x)=N_t^{\mathsf T}G_t\mathscr R_t(x)N_t\ .
 \label{eq:compressed-map}
\end{align}
This is a congruence of a positive physical moment block, multiplied by
$s$. For joint blocks, the same construction uses
$\operatorname{diag}(L_t,L_t)$ in \cref{eq:joint-linear} and a single
joint frame $N_t$. The off-diagonal entry is retained before compression.
The explicit metric-adjoint convention is stated in \cref{app:gram}.

No completeness property of the frames is needed for this positivity
implication. Any rational frame defines a valid positive compression.
Their selection and the equality families used to find them concern the
discovery of a certificate, not the coverage of the original matrices.
Completeness of the source coordinates in \cref{lem:source-complete}
ensures that a functional identity checked there holds for every
physical averaged moment.

\section{A rational Gram certificate for the three-copy endpoint}
\label{sec:certificate}

The finite certificate consists of rational symmetric matrices $Z_t$,
rational congruences $R_t$, and a rational equality multiplier $\eta$.
The frames, embeddings, source words, and equality row labels specify the
linear maps appearing with these multipliers.
The counts relevant to the proof are collected in \cref{tab:certificate}.

\begin{table}[ht]
\centering
\begin{tabular}{lr}
\toprule
Certificate quantity & Value\\
\midrule
Source dimensions $(31,22,13)$ & $(667,1407,667)$\\
Total source coordinates & $2741$\\
Equality matrix size & $2912\times2741$\\
Selected equality rows & $272$\\
Positive Gram multipliers & $1239$\\
Multipliers on joint blocks & $111$\\
Largest Gram order & $43$\\
Independent symmetric Gram entries & $18753$\\
Rows in the positivity checks & $4231$\\
\bottomrule
\end{tabular}
\caption{Exact certificate dimensions. These counts describe the supplied
construction and are not asserted to be minimal.}
\label{tab:certificate}
\end{table}

\begin{proposition}[Finite rational certificate]
\label{prop:certificate}
The supplied data satisfy the functional identity
\begin{align}
 c^{\mathsf T}x
 =\sum_{t=1}^{1239}\Tr\bigl[Z_t\widehat M_t(x)\bigr]
       +\eta^{\mathsf T}E_0x
 \qquad\text{for every }x\in\R^{2741}\ .
 \label{eq:exact-d4-three-certificate}
\end{align}
For each multiplier, $Z_t=Z_t^{\mathsf T}\in M_{k_t}(\Q)$ and its
saved rational congruence satisfies, for $1\leq i\leq k_t$,
\begin{align}
 (R_tZ_tR_t^{\mathsf T})_{ii}
 -\sum_{j\ne i}|(R_tZ_tR_t^{\mathsf T})_{ij}|
 >\frac{9999992}{10000000}\ .
 \label{eq:exact-congruence-positive}
\end{align}
In particular every $Z_t$ is positive definite.
\end{proposition}

\begin{proof}[Exact verification]
The source generators, literal module matrices, and frames determine
rational coefficient maps. They are rebuilt using integer contractions
with explicit bounds excluding overflow, as described in \cref{app:gram}.
For each standard coordinate vector $e_k\in\R^{2741}$, the verifier
computes the right side of \cref{eq:exact-d4-three-certificate} over
$\Q$ and checks equality to $c_k$. All 2741 differences are exactly zero.
Linearity proves the identity for arbitrary real coordinates.

The verifier reads the rational entries of $Z_t$ and $R_t$, checks their
dimensions and symmetry, and evaluates the left side of
\cref{eq:exact-congruence-positive} over $\Q$.
All 4231 inequalities satisfy the stated rational lower bound.
A real symmetric matrix with positive diagonal-dominance margins is
positive definite. Indeed, bounding each cross term in its quadratic
form by $2|v_i v_j|\leq v_i^2+v_j^2$ gives a strictly positive lower bound
for every nonzero vector.
Thus $R_tZ_tR_t^{\mathsf T}\succ0$.
This forces $R_t$ to be nonsingular, and congruence implies $Z_t\succ0$.
The saved inputs, the verification procedure, and the resulting records
are specified in \cref{app:verification}.
\end{proof}

\begin{proof}[Proof of \cref{thm:d4-three-certificate}]
For $C\ne0$, first normalize $\norm C_F^2=1/2$.
The lifting and common averaging in \cref{sec:matched-four} produce
physical moments. By \cref{lem:source-complete}, they have source
coordinates $x$ satisfying $Ex=0$ and $c^{\mathsf T}x=q_3(C)$.
Every $\widehat M_t(x)$ is positive semidefinite by
\cref{eq:literal-compression,eq:compressed-map}, with the same statement
for the joint compressions.
\Cref{prop:certificate} therefore gives
\begin{align}
 q_3(C)=\sum_{t=1}^{1239}
 \Tr\!\left[Z_t^{1/2}\widehat M_t(x)Z_t^{1/2}\right]\geq0\ .
\end{align}
Homogeneity restores an arbitrary norm, and $C=0$ is immediate.
Equation~\eqref{eq:distillation-form} and the finite-copy distillation
criterion complete the proof.
\end{proof}

The positivity of the stored multipliers is strict in their compressed
coordinates. The resulting physical inequality is not strict because
some physical moment compressions vanish, including those arising from
\cref{eq:product-zero}.
The proof uses no floating-point eigenvalue test, numerical rank decision,
or optimization stopping tolerance.
Numerical optimization was used to discover the data, after which an
exact affine correction and exact positive-congruence checks produced
the certificate. Details of that distinction are given in
\cref{app:discovery}.

\section{Discussion}
\label{sec:discussion}

The three-copy endpoint at local dimension four is controlled by an
exact positive identity on matched moments of arbitrary rank-two
coefficient matrices. The joint block relates the three moment families
through off-diagonal entries, while the conditional auxiliary-qubit
marginals encode the isometric support of the original matrix.
The dimension-four alternating component is retained throughout.
These ingredients permit a certificate that applies without structural
assumptions on the input matrix.

The result establishes the full three-copy threshold for the two-ququart
Werner family. It also applies to smaller local coefficient spaces that
embed isometrically in each $\C^4$ factor, since such embeddings preserve
all the partial-trace norms in \cref{eq:q3}.
For the NPT interval $-1/2\leq\alpha<-1/4$, collective processing of
three copies does not enable distillation.
The proof does not control four or arbitrarily many resource copies.
NPT bound entanglement requires undistillability at every finite copy
number and is therefore not established here.

Extending the certificate to larger copy numbers requires joint
control of the positive compressions and the matched-moment equalities.
Whether these constraints admit a positive recursion remains open.

\appendix

\section{Permutation conventions and homogeneous equalities}
\label{app:equalities}

\subsection{Diagram conventions}

All permutation lists in the data are lexicographically ordered.
A stored permutation is a tuple $(p(0),p(1),p(2),p(3))$, with labels
starting at zero. A stored word is a tuple of four indices into this
24-element list, in the order of three physical sites followed by $Q$.
The permutation matrices obey $V_pV_q=V_{p\circ q}$ and
$V_p^{\mathsf T}=V_{p^{-1}}$.
The maps in \cref{eq:joint-maps} therefore replace each local word entry
according to
\begin{align}
 p\longmapsto\tau p\tau^{-1}
 \quad\text{or}\quad p\longmapsto p\pi
 \label{eq:joint-word-maps}
\end{align}
before taking the indicated partial transpose.
In the stored labels, $\tau=(0,1,3,2)$ and $\pi=(2,3,1,0)$.
For a partially transposed diagram, permutation multiplication is
performed before transposition. The matrices representing partially
transposed diagrams are not assumed to obey the original permutation
group multiplication law.

The real local permutation span is denoted $\mathfrak C_d$.
For $d=4$, all 24 diagrams are independent because their actions on
$|0,1,2,3\rangle$ give distinct basis vectors.
For $d=2$, an independent list of permutation indices is
\begin{align}
 \mathcal I_2=(0,3,4,7,8,9,11,12,13,15,16,18,20,23)\ .
 \label{eq:auxiliary-basis}
\end{align}
Its Gram matrix is nonsingular, and all 24 literal matrices lie in its
span. Both facts are checked over $\Q$.
If $Q_{a,p}$ denotes the coordinate of $V_p^{(2)}$ in this list, then
\begin{align}
 V_p^{(2)}=\sum_{a=1}^{14}Q_{a,p}V_{\mathcal I_2(a)}^{(2)}\ .
 \label{eq:auxiliary-quotient}
\end{align}
The second implementation obtains $Q$ directly from the $16\times16$
permutation matrices and verifies the reconstructed matrices entrywise.

\subsection{Matching and conditional maps}

The equality maps can be specified before any coordinate reduction.
Let $\mathcal T_r$ trace the entire replica $\mathcal V_r$ and retain
the other replicas in increasing order. Let $\mathcal R$ on
$\mathcal T_2Y_{13}$ put its original labels in the order $(4,1,3)$.
The two matching equations are
\begin{align}
 \mathcal T_3Y_{31}-\mathcal T_3Y_{22}&=0,\nonumber\\
 \mathcal T_2Y_{22}-\mathcal R\mathcal T_2Y_{13}\mathcal R^{\mathsf T}&=0\ .
 \label{eq:matching-equalities}
\end{align}
For a selected $a$ replica, the conditional map is
\begin{align}
 \mathcal S_r(Y)=\Tr_{H_r}Y
 -\iota_r\!\left(\frac{I_{Q_r}}2\otimes\mathcal T_rY\right)\ ,
 \label{eq:conditional-map}
\end{align}
where $\iota_r$ restores the original positions of the remaining factors.
The three conditional equations are
$\mathcal S_1(Y_{31})=0$, $\mathcal S_1(Y_{22})=0$, and
$\mathcal S_4(Y_{13})=0$.

After tracing the physical factors of one replica, each remaining
physical site has three replicas. Its six permutation diagrams are
independent at $d=4$.
Averaging over physical-site permutations leaves
$\binom{6+3-1}{3}=56$ orbit coordinates.
The untraced four-replica auxiliary algebra has dimension 14, giving
$56\cdot14=784$ output coordinates for each conditional map.
In a matching map the auxiliary system also has three replicas.
Its permutation span has dimension five, because the three-fold
antisymmetrizer vanishes at $d=2$. Thus each matching map has
$56\cdot5=280$ coordinates.
The total output size is $3\cdot784+2\cdot280=2912$.

The stored row order consists of the three conditional maps for
$(31,22,13)$, followed by the $aab$ and $abb$ matching maps.
Within each output, physical permutation triples are in increasing
lexicographic order after sorting, followed by the auxiliary index.
Coefficients refer to orbit averages, rather than orbit sums.
Expanding the five maps on each source basis generator and clearing the
common denominator $1152$ gives the integer matrix $E$.
The selected row labels in the certificate also start at zero.
The multiplier $\eta$ is paired with these integer rows, so no additional
factor $1152$ belongs in \cref{eq:exact-d4-three-certificate}.

The model uses these homogeneous equalities. Although physical moments
have unit trace, an affine trace constraint is not used in the
certificate identity. Multiplication of all moments by the positive
constant $s$ in \cref{eq:matched-source-scale} is therefore compatible
with every equality in the proof.
The source and equality data are reconstructed from literal partial
traces in the independent verifier, including all orderings in
\cref{eq:matching-equalities}.

\section{Completeness of the source basis}
\label{app:source}

\subsection{An exact projection-rank calculation}

Before applying the identical-factor projectors, the physical-site
invariant part of the commutant is
$\Sym^3(\mathfrak C_4)\otimes\mathfrak C_2$.
The three operations in \cref{eq:source-projection} commute after
averaging over the identical-factor subgroups. The transpose operation
commutes with the sandwich projector because $\Pi_f$ is real symmetric.
Physical-site permutations commute with simultaneous replica permutations.
Consequently $\mathcal P_f$ is idempotent on the relevant commutant.

For $h,k\in R_f$ and $\epsilon\in\{0,1\}$, the local linear map is
\begin{align}
 K_{h,k,\epsilon}^{(d)}(X)=V_h^{(d)}X^{\mathsf T\epsilon}V_k^{(d)}\ ,
 \label{eq:local-linear-action}
\end{align}
where $X^{\mathsf T0}=X$ and $X^{\mathsf T1}=X^{\mathsf T}$.
Its traces on the two local permutation spans determine the global
projection trace. With
$t_j=\Tr[(K_{h,k,\epsilon}^{(4)})^j]$ and
$u=\Tr[K_{h,k,\epsilon}^{(2)}]$, the symmetric-cube identity is
\begin{align}
 \Tr\bigl[\Sym^3(K_{h,k,\epsilon}^{(4)})\bigr]
 =\frac{t_1^3+3t_1t_2+2t_3}{6}\ .
 \label{eq:symmetric-cube-trace}
\end{align}
To derive it, insert the average of the six factor permutations into
$K^{\otimes3}$. The identity permutation contributes $t_1^3$, the three
transpositions contribute $t_1t_2$, and the two three-cycles contribute
$t_3$.
Therefore
\begin{align}
 \dim\ran\mathcal P_f
 =\frac1{2|R_f|^2}\sum_{h,k\in R_f}\sum_{\epsilon=0}^1
       \frac{t_1^3+3t_1t_2+2t_3}{6}\,u\ .
 \label{eq:source-rank-formula}
\end{align}
This calculation needs only literal permutation actions and the auxiliary
quotient in \cref{eq:auxiliary-quotient}.
At $d=4$, $K$ sends $V_p$ to $V_{hpk}$ when $\epsilon=0$ and to
$V_{hp^{-1}k}$ when $\epsilon=1$.
Its power traces count fixed labels in the 24-element basis.
At $d=2$, the same images are reduced in the 14-element basis before
taking the trace.
The exact sums are given in \cref{tab:source-traces}.

\begin{table}[ht]
\centering
\begin{tabular}{crrrr}
\toprule
Family & $\epsilon=0$ sum & $\epsilon=1$ sum & $2|R_f|^2$ & Rank\\
\midrule
$31$ & $37584$ & $10440$ & $72$ & $667$\\
$22$ & $37984$ & $7040$ & $32$ & $1407$\\
$13$ & $37584$ & $10440$ & $72$ & $667$\\
\bottomrule
\end{tabular}
\caption{The two sums include the factor $1/6$ from
\cref{eq:symmetric-cube-trace}. Adding them and dividing by
$2|R_f|^2$ gives the exact source dimension.}
\label{tab:source-traces}
\end{table}

\subsection{Independent selected generators}

For each word in the saved basis list, the verifier reconstructs
$\mathcal P_f(D_{\boldsymbol p})$ by enumerating both sandwich subgroup
elements, all six physical-site permutations, and the transpose choice.
The common diagram denominator is $12|R_f|^2$, hence 432 for $31,13$
and 192 for $22$.
The resulting full diagram coefficient matrix is compared exactly with
the stored source table.

For the independence check, the auxiliary diagrams are reduced through
\cref{eq:auxiliary-quotient} and identical physical orbits are collected.
If an orbit coefficient is collected as a total weight, its basis element
is the corresponding orbit average. This convention fixes the factors
for orbits containing repeated permutations.
After clearing denominators in each column, the saved pivot rows select
a square integer matrix. Its determinant is evaluated modulo $1000033$.
The nonzero values quoted in \cref{lem:source-complete} prove that the
integer determinants are nonzero.
Only this lower-bound implication is used from modular arithmetic.
The matching upper bounds come from the rational idempotent trace,
not from a numerical or modular rank estimate alone.

Finally, the objective can be reconstructed independently of any moment
sample. For each generator,
$c_{31,k}=s^{-1}\Tr(OB_{31,k})$ is evaluated using literal local matrix
traces, while the other two families have zero objective coefficients.
For a local permutation, $\Tr V_p^{(d)}=d^{\#\mathrm{cycles}(p)}$.
Insertion of a swap replaces $p$ by its product with that transposition.
These exact traces verify the normalization as well as the sign of every
term in \cref{eq:q3}.

\section{Local embeddings and reconstruction of the Gram maps}
\label{app:gram}

\subsection{Physical meaning of the saved representations}

The retained source cuts, stated with replica labels starting at one,
are listed in \cref{tab:cuts}. Every cut transposes an entire replica of
$\mathcal V$, including its auxiliary qubit. Joint blocks use the
single transpose at replica 4 that occurs in \cref{eq:joint-maps}.
A transpose at only a physical subfactor would not have the positivity
justification used in this proof.

\begin{table}[ht]
\centering
\begin{tabular}{cl}
\toprule
Family & Complete-replica cuts\\
\midrule
$31,13$ & $\varnothing,\{1\},\{4\},\{1,2\}$\\
$22$ & $\varnothing,\{1\},\{3\},\{1,2\},\{1,3\}$\\
Joint & $\{4\}$ after the word maps in \cref{eq:joint-word-maps}\\
\bottomrule
\end{tabular}
\caption{Positive maps used in the certificate. Selecting a subset of
valid positive maps does not restrict the original physical inputs.}
\label{tab:cuts}
\end{table}

For each cut type, the local module dimensions are
\begin{align}
 d=4\quad&
 \begin{cases}
 1,1,2,3,3 & |T|=0,1,\\
 1,1,1,1,2,4 & |T|=2,
 \end{cases}\nonumber\\
 d=2\quad&1,2,3\qquad(|T|=0,1,2)\ .
 \label{eq:local-module-sizes}
\end{align}
The embeddings for the different positions of a cut are included
explicitly. Their squared dimensions sum to 24 at $d=4$ and 14 at $d=2$.
The positivity argument uses the literal full-column-rank embeddings,
so it does not require a classification or completeness assertion for
the local module list.

All entries of $J_\lambda$, $G_\lambda$, and $\rho_\lambda(p)$ are
rational. The verifier constructs $(V_p^{(d)})^{\Gamma_T}$ from the
identity tensor by permuting output indices and then exchanging the
row and column indices at each transposed replica.
It checks \cref{eq:local-intertwiner} and full column rank of $J_\lambda$.
Thus $G_\lambda\succ0$ follows from a literal Gram matrix, without an
eigenvalue tolerance.

In nonorthonormal coordinates the adjoint relation is
\begin{align}
 G_t\mathscr R_t(X^\dagger)
       =\mathscr R_t(X)^\dagger G_t\ .
 \label{eq:metric-adjoint}
\end{align}
For real averaged source coordinates, let $A_t(x),C_t(x),D_t(x)$ be the
represented top diagonal, cross, and bottom diagonal entries of a joint
block. Before the saved joint frame is applied, the symmetric compression
is
\begin{align}
 H_t(x)=
 \begin{pmatrix}
 G_tA_t(x)&G_tC_t(x)\\
 C_t(x)^{\mathsf T}G_t&G_tD_t(x)
 \end{pmatrix}\ .
 \label{eq:joint-metric-block}
\end{align}
It equals a congruence of the literal physical joint block by
$\operatorname{diag}(L_t,L_t)$.
The positive map used in the certificate is
$\widehat M_t(x)=N_t^{\mathsf T}H_t(x)N_t$.
In particular the lower cross entry is the transpose of $G_tC_t(x)$,
not a bare transpose of a nonorthonormal representation followed by an
unjustified metric placement.

Before the final frames, the selected construction has 1722 parent
blocks, of which 702 involve an additional dimension-four local module.
The retained maps comprise 335 source blocks for each of $31,13$, 458
for $22$, and 111 joint blocks. Different frames may be associated with
one parent block. Of the 1239 retained maps, 522 involve an additional
dimension-four local module.
The archive gives each parent, cut, local module tuple, frame, and order
explicitly. These labels are checked when a multiplier is read.

\subsection{Rational coefficient reconstruction}

For a Gram matrix of order $k_t$, the stored symmetric coordinates are
$Z_{ij}$ with $j=1,\ldots,k_t$ and $i=j,\ldots,k_t$, in that order.
Let $E^{[ij]}$ be $E_{ii}$ for $i=j$ and $E_{ij}+E_{ji}$ otherwise.
The coefficient map is determined by
\begin{align}
 (L_t)_{k,[ij]}
 =\Tr\!\left[E^{[ij]}\widehat M_t(e_k)\right]\ .
 \label{eq:gram-coefficients}
\end{align}
For an off-diagonal coordinate this includes both symmetric entries.
No $\sqrt2$ factor is used in these exact stored coordinates.

An explicit reconstruction uses $W=N_tE^{[ij]}N_t^{\mathsf T}$.
For a source block it contracts $W$ against the tensor products of the
local matrices $G_\lambda\rho_\lambda(p)$, followed by the rational
source diagram table.
For a joint block, partition $W$ into its two-by-two blocks. The pairing
with \cref{eq:joint-metric-block} is
\begin{align}
 \Tr(W H_t)
 =\Tr(W_{11}G_tA_t)
  +2\Tr(W_{21}G_tC_t)
  +\Tr(W_{22}G_tD_t)\ .
 \label{eq:joint-coefficient-pairing}
\end{align}
This accounts for the factor two in every joint cross contribution.
The word maps in \cref{eq:joint-word-maps} are applied locally before the
partially transposed representations are selected.

All input rational matrices are first expressed as integer matrices with
common positive denominators. Tensor contractions then use integer
arithmetic. Before a product $AB$, an integer bound
\begin{align}
 |(AB)_{ij}|
 \leq \left(\sum_\ell|A_{i\ell}|\right)
           \max_{\ell,j}|B_{\ell j}|
 \label{eq:integer-bound}
\end{align}
controls every intermediate sum as well as the final entry.
The second implementation checks bounds with arbitrary-precision Python
integers before using fixed-width integer products. All accepted bounds
are below $2^{62}$.
Sparse source contractions are bounded in the same way.
Denominators are kept separately and common factors are canceled exactly.
Thus the independently reconstructed coefficient maps contain no
coefficients inferred by rounding a floating-point result.

The final identity check evaluates
\begin{align}
 r_k=c_k-\sum_t\Tr\!\left[Z_t\widehat M_t(e_k)\right]
                  -\eta^{\mathsf T}E_0e_k
 \label{eq:residual-check}
\end{align}
for all $k=1,\ldots,2741$ in rational arithmetic.
Every $r_k$ is exactly zero. There is no omitted-coordinate or
least-squares tolerance in this check.
Together with \cref{lem:source-complete}, this verifies the identity on
all averaged physical inputs.

\section{Discovery of the certificate and exact reconstruction}
\label{app:discovery}

\subsection{Equality states and multiplier ranges}

Known zero states help identify useful multiplier ranges in a degenerate
semidefinite problem. If an identity with positive multipliers holds at a
physical zero $x_*$, then each term
$\Tr[Z_t\widehat M_t(x_*)]$ is nonnegative and their sum is zero.
It follows that
\begin{align}
 \ran Z_t\subseteq\ker\widehat M_t(x_*)\ .
 \label{eq:face-condition}
\end{align}
Indeed, zero trace of a product of positive matrices gives
$\widehat M_t(x_*)^{1/2}Z_t^{1/2}=0$.
Intersecting such kernel constraints implements the relevant part of
facial reduction~\cite{PermenterParrilo}.
The product zeros in \cref{eq:product-zero} already impose restrictions,
but additional nonnormal zero families constrain further cross entries.

A useful family lies on embedded qubit supports.
Let $J=\begin{pmatrix}0&1\\-1&0\end{pmatrix}$ and
$\Theta=J^{\otimes3}\mathsf K$, where $\mathsf K$ is complex conjugation.
The antiunitary obeys $\Theta^2=-I$.
For unit $u,v\in(\C^2)^{\otimes3}$ satisfying
$\ip u{\Theta v}=0$, define
\begin{align}
 C_{u,v}=\frac12\bigl(|u\rangle\langle v|
             +|\Theta v\rangle\langle\Theta u|\bigr)\ .
 \label{eq:crossed-antistate}
\end{align}
The two pairs $(u,\Theta v)$ and $(v,\Theta u)$ are orthonormal.
Thus $C_{u,v}$ has rank two and squared Frobenius norm $1/2$.
Its zero objective follows from the linear involution
$\mathcal S(C)=J^{\otimes3}C^{\mathsf T}(J^{\otimes3})^\dagger$.
This involution exchanges the two summands and acts by $(-1)^w$ on a
Pauli string of weight $w$. The weight-three coefficients of $C_{u,v}$
therefore vanish.

To translate that fact into $q_3$, the maps
$\mathcal D_i(C)=C-(I_{2,i}/2)\otimes\Tr_iC$ are commuting
orthogonal projections on the embedded qubit matrix space.
Their product selects precisely weight-three Pauli strings, and expansion
gives
\begin{align}
 q_3(C)=\ip C{\mathcal D_1\mathcal D_2\mathcal D_3(C)}
       =\norm{\mathcal D_1\mathcal D_2\mathcal D_3(C)}_F^2\ .
\end{align}
Hence $q_3(C_{u,v})=0$.
Local isometric embedding in $\C^4$ preserves every partial-trace norm,
so the same zero is valid for the target dimension.
The balanced choice $v=u$ gives the anti-state family discussed in
Ref.~\cite{WZ}. Allowing distinct $u,v$ also gives nonnormal members.

For example, a real seed and its paired vector are
\begin{align}
 u=\frac15(1,2,-1,3,2,-2,1,1)^{\mathsf T},
 \qquad v=(\sigma_z\otimes\sigma_z\otimes I_2)u\ .
 \label{eq:crossed-seed}
\end{align}
They satisfy the required orthogonality.
For $C_{u,v}=C_0/50$, the sums of squared partial-trace norms of the
integer matrix $C_0$ at subset sizes $0,1,2,3$ are
$(1250,2742,726,484)$.
The identity $8(1250)-4(2742)+2(726)-484=0$ gives its exact zero.
It is nonnormal because
\begin{align}
 \norm{C_{u,v}C_{u,v}^\dagger-C_{u,v}^\dagger C_{u,v}}_F^2
 =\frac{1260000}{50^4}>0\ .
\end{align}
These families were used to select the frames supplied with the
certificate. Rational averaging of parameter families can also use
unitary permutation integration~\cite{CollinsSniady}.
Their role is to guide the search. The final proof needs only that each
saved frame is a valid congruence and that the complete functional
identity and multiplier positivity hold.

\subsection{From a numerical candidate to a rational identity}

The numerical search optimized complete Gram matrices in the retained
frames. A candidate with a prescribed positive margin in its numerical
coordinates was converted to rational Gram factors by dyadic rounding.
Factor products give initially positive semidefinite rational matrices.
They do not automatically satisfy the desired coefficient identity.

An exact affine correction was then computed together with the equality
multipliers. After clearing denominators in the coefficient map, a chosen
$2634\times2634$ integer minor was solved over $\Q$.
The resulting correction was checked against all 2741 original source
coordinates, including the 107 coordinates outside that minor.
Thus a successful solve of a selected subsystem was not treated as a
proof of the full identity.

An affine correction can destroy positivity. After the correction, a
separate verification therefore tested every final Gram matrix by its
saved rational congruence. The lower bound in
\cref{eq:exact-congruence-positive} applies to these final matrices.
Its value is a margin for the transformed Gram matrices, not a positive
lower bound for $q_3(C)$, which has exact zeros.
The final certificate can be checked without rerunning any optimization,
choosing a numerical rank threshold, or reproducing this correction.

\section{Certificate data and reproducible verification}
\label{app:verification}

\subsection{Data needed to reproduce the proof}

The companion repository\footnote{\url{https://github.com/GLiuax/d4-three-copy-werner}} contains
the manuscript source, the exact data, and the independent verifier.
It preserves the relative directory \nolinkurl{research/matched_d4}.
Within that directory, \nolinkurl{source.json} gives the selected words,
normalization, objective, and source-basis information.
The three \nolinkurl{diagrams_*.npz} files store integer source diagram
tables with their denominators specified in the source file.
The file \nolinkurl{equality_integer.npz} contains $E$.
Local embeddings and diagram images are in
\nolinkurl{aligned_modules.json}, and parent block labels and replica
cuts are in \nolinkurl{positive_blocks.json}.

The subdirectory \nolinkurl{full_gram_crossed} contains
\nolinkurl{exact_frames.json} and \nolinkurl{frames.json}, which specify
the rational frames and their block labels.
Its \nolinkurl{rational_maps} directory contains the rational Gram
coefficient maps, represented as integer matrices with scalar
denominators.
The \nolinkurl{rational_certificate} directory contains files
\nolinkurl{0000.json} through \nolinkurl{1238.json}, each recording a
multiplier $Z_t$, its congruence $R_t$, and its labels.
The file \nolinkurl{equality.json} records the selected integer row
indices and the rational entries of $\eta$.
Rational entries are serialized as integer or fraction strings.

The original input manifest records 2485 hashes covering the multipliers
and their coordinate inputs. The independent verifier checks this
pre-existing manifest without replacing it.
The source diagram tables are additionally reconstructed from their
literal defining words.
The companion repository includes a separate package manifest for the
manuscript, verifier, and all distributed files.
Hashes identify an exact set of inputs. Their agreement is an integrity
check and is not substituted for any algebraic verification.

The dimension-four source-coordinate construction used permutation
conventions and auxiliary quotient data from the MIT-licensed code
accompanying Ref.~\cite{Candidate}. The dimension-four literal embeddings,
additional source directions, frames, and present certificate were
constructed in this project.
The independent audit program imports no construction program or earlier
certificate verifier. It uses literal permutations, the serialized
rational inputs, and general-purpose arithmetic libraries.
The companion repository retains the upstream attribution and license
information for the reused material.

\subsection{Verification stages}

The independent program is
\nolinkurl{research/d4_independent_certificate_audit.py}.
The five stages can be run from the repository root.
\begin{verbatim}
cd research
audit=d4_independent_certificate_audit.py
python "$audit" reduction
python "$audit" foundation
python "$audit" rank
python "$audit" maps --workers 3
python "$audit" certificate
\end{verbatim}
The stages perform the following checks.
\begin{enumerate}
\item The reduction stage compares formal matrix-entry monomials for all
three joint identities. It also performs numerical diagnostics of the
swap and vectorization conventions on general complex inputs, including
rank one and unequal singular values. The numerical diagnostics are not
premises of the universal proof, which is the index argument in
\cref{sec:matched-four}.
\item The foundation stage checks the original input hashes, all 50
literal module embeddings and 1200 intertwining equations, the source
projection dimensions, and nonzero basis minors. It reconstructs every
source generator and every equality and objective column using literal
partial traces and matrix traces.
\item The rank stage verifies over $\Q$ that both $E$ and the selected
row submatrix have rank 272. For the positivity implication, validity of
the selected rows alone would suffice.
\item The maps stage rebuilds all 1239 rational Gram maps using the
integer contractions described in \cref{app:gram}. It checks exact
agreement with the saved maps and writes the independently rebuilt maps
to a separate audit directory.
\item The certificate stage uses the rebuilt maps and objective to check
all 2741 rational identities. It checks symmetry and the strict rational
diagonal-dominance inequalities for all 1239 multipliers, covering 4231
rows in total.
\end{enumerate}

The recorded rational run used NumPy 2.4.4, SciPy 1.17.1, SymPy 1.14.0,
and python-flint 0.9.0.
The coefficient identity and final positivity checks use exact rational
arithmetic. Integer operations used for the map reconstruction are
bounded as in \cref{eq:integer-bound}.
The earlier verifier additionally checked positivity by rigorous Arb
interval arithmetic at 256-bit precision, but the rational checks alone
suffice for the certificate.

All reported checks passed on 14 October 2026.
The records are stored in
\nolinkurl{research/matched_d4/full_gram_crossed/audit_2026-10-14}.
The two implementations were developed within the same project.
Their agreement supplies implementation diversity and does not constitute
external peer review or formal verification in a proof assistant.
The mathematical reduction, source-space completeness, and the finite
rational checks together constitute the computer-assisted proof.

\begin{thebibliography}{99}

\bibitem{HHH}
M. Horodecki, P. Horodecki, and R. Horodecki,
Mixed-state entanglement and distillation: Is there a ``bound''
entanglement in nature?,
Phys. Rev. Lett. \textbf{80}, 5239--5242 (1998).
\href{https://doi.org/10.1103/PhysRevLett.80.5239}{doi 10.1103/PhysRevLett.80.5239}.

\bibitem{Werner}
R. F. Werner,
Quantum states with Einstein--Podolsky--Rosen correlations admitting
a hidden-variable model,
Phys. Rev. A \textbf{40}, 4277--4281 (1989).
\href{https://doi.org/10.1103/PhysRevA.40.4277}{doi 10.1103/PhysRevA.40.4277}.

\bibitem{HHWerner}
M. Horodecki and P. Horodecki,
Reduction criterion of separability and limits for a class of
distillation protocols,
Phys. Rev. A \textbf{59}, 4206--4216 (1999).
\href{https://doi.org/10.1103/PhysRevA.59.4206}{doi 10.1103/PhysRevA.59.4206}.

\bibitem{DCLB}
W. D\"ur, J. I. Cirac, M. Lewenstein, and D. Bru\ss,
Distillability and partial transposition in bipartite systems,
Phys. Rev. A \textbf{61}, 062313 (2000).
\href{https://doi.org/10.1103/PhysRevA.61.062313}{doi 10.1103/PhysRevA.61.062313}.

\bibitem{DSS}
D. P. DiVincenzo, P. W. Shor, J. A. Smolin, B. M. Terhal, and A. V. Thapliyal,
Evidence for bound entangled states with negative partial transpose,
Phys. Rev. A \textbf{61}, 062312 (2000).
\href{https://doi.org/10.1103/PhysRevA.61.062312}{doi 10.1103/PhysRevA.61.062312}.

\bibitem{Five}
P. Horodecki, L. Rudnicki, and K. \.Zyczkowski,
Five open problems in quantum information theory,
PRX Quantum \textbf{3}, 010101 (2022).
\href{https://doi.org/10.1103/PRXQuantum.3.010101}{doi 10.1103/PRXQuantum.3.010101}.

\bibitem{PPHH}
\L{}. Pankowski, M. Piani, M. Horodecki, and P. Horodecki,
A few steps more towards NPT bound entanglement,
IEEE Trans. Inf. Theory \textbf{56}, 4085--4100 (2010).
\href{https://doi.org/10.1109/TIT.2010.2050810}{doi 10.1109/TIT.2010.2050810}.

\bibitem{CostaRico}
P. Costa Rico,
New partial trace inequalities and distillability of Werner states,
Lett. Math. Phys. \textbf{115}, 47 (2025).
\href{https://doi.org/10.1007/s11005-025-01935-y}{doi 10.1007/s11005-025-01935-y}.

\bibitem{SC}
Z. Song and L. Chen,
A partial-trace matrix inequality and Werner-state distillability,
\href{https://arxiv.org/abs/2607.23416}{arXiv 2607.23416} (2026).

\bibitem{FHPV}
T. C. Fraser, F. Huber, B. Pozsgay, and I. Vona,
On the two-copy distillability of Werner states and a new partial trace inequality,
\href{https://arxiv.org/abs/2607.24309}{arXiv 2607.24309} (2026).

\bibitem{FGP}
J. Fu, L. Gao, and S.-J. Park,
A solution to 2-copy distillability of Werner states,
\href{https://arxiv.org/abs/2607.21367}{arXiv 2607.21367} (2026).

\bibitem{BGH}
K. Bharti, R. Gajjala, and T. Haug,
Two-copy nondistillability of Werner states:
Sharp partial-trace inequalities and finite-copy extensions,
\href{https://arxiv.org/abs/2607.24479}{arXiv 2607.24479} (2026).

\bibitem{WZ}
T. Wu and Q. Zou,
Sharp Pl\"ucker geometry for three-copy Werner distillation,
\href{https://arxiv.org/abs/2608.02647}{arXiv 2608.02647} (2026).

\bibitem{VD}
R. O. Vianna and A. C. Doherty,
Distillability of Werner states using entanglement witnesses and robust
semidefinite programs,
Phys. Rev. A \textbf{74}, 052306 (2006).
\href{https://doi.org/10.1103/PhysRevA.74.052306}{doi 10.1103/PhysRevA.74.052306}.

\bibitem{Candidate}
Anonymous,
Three-copy undistillability of the boundary qutrit Werner state,
version 1.1.0-candidate (4 October 2026), candidate manuscript and certificate.
\href{https://doi.org/10.5281/zenodo.23134454}{doi 10.5281/zenodo.23134454}.

\bibitem{PermenterParrilo}
F. Permenter and P. A. Parrilo,
Partial facial reduction: Simplified, equivalent SDPs via approximations
of the PSD cone,
\href{https://arxiv.org/abs/1408.4685}{arXiv 1408.4685} (2014).

\bibitem{CollinsSniady}
B. Collins and P. \'Sniady,
Integration with respect to the Haar measure on unitary, orthogonal and
symplectic group,
Commun. Math. Phys. \textbf{264}, 773--795 (2006).
\href{https://doi.org/10.1007/s00220-006-1554-3}{doi 10.1007/s00220-006-1554-3}.

\end{thebibliography}
\end{document}
